\section{Introduction}\label{introduction}

\subsection{Clinical Context and Motivation}\label{clinical-context-and-motivation}

Autism Spectrum Disorder (ASD) is a neurodevelopmental condition characterized by persistent deficits in social communication and restricted, repetitive behaviors \cite{ref1}. With global prevalence rising to 1 in 36 children \cite{ref2}, early diagnosis has become increasingly critical for enabling timely intervention \cite{ref3}. However, current diagnostic procedures---the Autism Diagnostic Observation Schedule (ADOS) and Autism Diagnostic Interview-Revised (ADI-R)---require 2-3 hours per assessment, specialized clinical training, and remain inaccessible in resource-limited settings \cite{ref4}.

The pressing need for accessible, objective screening tools has driven significant interest in machine learning (ML) approaches for ASD diagnosis \cite{ref5,ref6,ref7,ref8,ref9}. Recent studies have demonstrated promising results using ensemble methods \cite{ref10}, deep learning \cite{ref11}, and evolutionary algorithms \cite{ref12}. However, critical challenges persist that limit clinical translation.

\subsection{Limitations of Existing Approaches}\label{limitations-of-existing-approaches-1}

\subsubsection{Multicollinearity in Behavioral Data}

Screening instruments like the Autism Spectrum Quotient (AQ-10) contain naturally correlated items that measure overlapping aspects of autism symptomatology \cite{ref13}. While multicollinearity is extensively studied in regression analysis \cite{ref14}, its role in classification-based medical feature selection remains underexplored. Standard approaches either ignore it entirely or apply generic thresholds (VIF\textless10) developed for non-clinical datasets \cite{ref15}.

\subsubsection{Interpretability Gap}

Clinical applications require transparent explanations that clinicians can evaluate and trust \cite{ref16}. While explainable AI (XAI) methods such as SHAP \cite{ref17} and LIME \cite{ref18} have been applied to autism diagnosis \cite{ref19,ref20,ref21}, few studies align ML-selected features with established clinical diagnostic frameworks such as the DSM-5 \cite{ref1}.

\subsubsection{Age Heterogeneity}

Autism presents differently across developmental stages, with social communication deficits manifesting distinctly in children versus adults \cite{ref22}. Most existing ML approaches treat all ages identically, missing developmentally specific diagnostic patterns \cite{ref23}.

\subsection{Research Contributions}\label{research-contributions-1}

To address these gaps, we present \textbf{EVO-XAI} (Explainable AI framework), an innovative approach that integrates:

\begin{enumerate}
\def\labelenumi{\arabic{enumi}.}
\item
  \textbf{Context-Specific VIF Optimization}: Systematic evaluation of VIF thresholds (5-13) reveals that autism screening requires different multicollinearity tolerance than regression-based benchmarks. While hyperparameter tuning is standard practice, prior ML work universally applied VIF\textless10 (inherited from regression standards) without questioning its appropriateness for behavioral symptom clustering. This work empirically demonstrates that VIF=13 is optimal for the autism domain, with the multi-seed verified average MCC reaching 0.816.
\item
  \textbf{DSM-5 Clinical Mapping}: First study explicitly aligning ML-selected features with DSM-5 diagnostic criteria, enabling clinical interpretability.
\item
  \textbf{Age-Stratified Validation}: Comprehensive evaluation across child, adolescent, and adult populations with developmental stage-specific feature selection.
\item
  \textbf{Robustness Guarantee}: Verification of model stability under significant noise perturbation (up to 20\%), addressing data quality concerns in clinical settings.
\end{enumerate}

\subsection{Paper Organization}\label{paper-organization}

Section 2 reviews related work. Section 3 details the proposed methodology. Section 4 presents experimental results. Section 5 discusses clinical implications and the VIF=13 phenomenon. Section 6 concludes with future directions.

\section{Related Work}\label{related-work}

\subsection{Machine Learning for Autism Diagnosis}\label{machine-learning-for-autism-diagnosis}

The application of ML to ASD screening has evolved significantly over the past decade. Early work by Wall et al.~\cite{ref24} demonstrated that ML could reduce ADOS administration time while maintaining diagnostic accuracy. Thabtah and colleagues \cite{ref5,ref25} developed the AQ-10 dataset and applied various supervised learning algorithms, establishing benchmark performance metrics.

Recent advances include ensemble methods achieving MCC values of 0.55-0.65 \cite{ref10,ref26}, deep learning approaches with neural attention mechanisms \cite{ref11,ref27}, and federated learning implementations for privacy-preserving diagnosis \cite{ref28}. A comprehensive 2024 review by Batsakis et al.~\cite{ref29} emphasized the integration of XAI to improve clinician trust. Qiang et al.~\cite{ref30} achieved 82.1\% classification accuracy using hierarchical brain network interactions from fMRI data. Notably, Rajagopalan et al.~\cite{ref31} reported a large-scale validation on the SPARK consortium (n=30,660), achieving AUROC 0.895 using logistic regression with XGBoost. However, this work did not address gaps that EVO-XAI targets, specifically, multicollinearity optimization, DSM-5 clinical mapping, or age-stratified performance. Interpretability remains a persistent challenge across these approaches.

\subsection{Explainable AI in Medical Diagnosis}\label{explainable-ai-in-medical-diagnosis}

The XAI field has produced several prominent techniques, including SHAP (Shapley Additive exPlanations)~\cite{ref17}, LIME (Local Interpretable Model-agnostic Explanations)~\cite{ref18}, and attention-based visualization~\cite{ref32}. SHAP has become the most widely used post-hoc explainer in clinical decision support systems~\cite{ref33}. In the context of autism diagnosis specifically, recent work has applied SHAP to interpret ensemble predictions~\cite{ref34}~and developed semantic web-enabled frameworks for scalable detection~\cite{ref20}. As highlighted by Rudin~\cite{ref16}, black-box models are inappropriate for high-stakes medical decisions, motivating our emphasis on interpretable feature selection and clinical alignment.

\subsection{Evolutionary Algorithms for Feature Selection}\label{evolutionary-algorithms-for-feature-selection}

Multi-objective evolutionary algorithms have proven effective for feature selection~\cite{ref35,ref36}~in medical applications~\cite{ref37,ref38}. NSGA-II~\cite{ref39}~provides Pareto-optimal solutions balancing accuracy and feature parsimony. Recent applications include COVID-19 prediction~\cite{ref37}~and breast cancer HER2 classification~\cite{ref38}, demonstrating NSGA-II's effectiveness in balancing feature minimization with classification accuracy. Differential Evolution~\cite{ref40}~and Particle Swarm Optimization~\cite{ref41}~offer alternative search strategies with different exploration-exploitation trade-offs. In autism research, bio-inspired and evolutionary approaches are increasingly recognized as a means of methodological innovation, with recent studies demonstrating their efficacy in feature selection and parameter optimization \cite{ref42,ref43,ref44}.

\subsection{Multicollinearity in Medical ML}\label{multicollinearity-in-medical-ml}

Multicollinearity, high correlation among predictor variables, is well-studied in statistical regression~\cite{ref14,ref45}. Recent research has demonstrated that VIF-based feature elimination significantly improves classification accuracy for Random Forest, SVM, and Gaussian Mixture Models~\cite{ref46}. However, standard practice recommends VIF \textless{} 10~\cite{ref15}, a threshold developed for regression coefficient estimation rather than for optimizing classification performance. Our work contributes the first systematic study of VIF threshold optimization for behavioral screening instruments, demonstrating that context-specific thresholds outperform universal standards.

\section{Materials and Methods}\label{materials-and-methods}

\subsection{Datasets}\label{datasets}

Three publicly available UCI datasets were employed \cite{ref47,ref48,ref49}. Table 1 presents the dataset characteristics.

\textbf{Table 1: Dataset Characteristics}

\begin{longtable}[]{@{}
  >{\raggedright\arraybackslash}p{(\linewidth - 6\tabcolsep) * \real{0.3066}}
  >{\raggedright\arraybackslash}p{(\linewidth - 6\tabcolsep) * \real{0.2324}}
  >{\raggedright\arraybackslash}p{(\linewidth - 6\tabcolsep) * \real{0.2287}}
  >{\raggedright\arraybackslash}p{(\linewidth - 6\tabcolsep) * \real{0.2324}}@{}}
\toprule\noalign{}
\begin{minipage}[b]{\linewidth}\raggedright
Characteristic
\end{minipage} & \begin{minipage}[b]{\linewidth}\raggedright
Child
\end{minipage} & \begin{minipage}[b]{\linewidth}\raggedright
Adolescent
\end{minipage} & \begin{minipage}[b]{\linewidth}\raggedright
Adult
\end{minipage} \\
\midrule\noalign{}
\endhead
\bottomrule\noalign{}
\endlastfoot
Sample Size & 292 & 104 & 704 \\
ASD Cases (\%) & 141 (48.3\%) & 63 (60.6\%) & 189 (26.8\%) \\
Age Range & 4-11 years & 12-16 years & 17+ years \\
Imbalance Ratio & 1.07:1 & 1.54:1 & 2.73:1 \\
\end{longtable}

Each dataset contained 20 features: 10 AQ-10 behavioral items, 9 demographic variables, and family history indicator. Fig. 1 displays the correlation structure of features in the child dataset, illustrating the inherent multicollinearity among AQ-10 items.

\textbf{Study Design Classification:} This is a secondary analysis of pre-existing diagnosed cohorts, not prospective clinical validation. The UCI autism datasets were developed by Thabtah \cite{ref47,ref48,ref49} using questionnaire responses from clinically confirmed autism cases and controls. True prospective clinical validation would require: (1) prospective recruitment of suspected autism cases, (2) blinded structured interviews (e.g., ADOS-2), and (3) comparison of model predictions to independent clinician diagnosis.

\begin{figure}[tbp]
\centering
\includegraphics[width=0.88\linewidth]{Figure_1.png}
\caption{Feature correlation matrix for the child dataset, showing moderate-to-high correlations among AQ-10 behavioral items (A1-A10).}
\label{fig:correlation-child}
\end{figure}

\subsection{Three-Stage Feature Selection Pipeline}\label{three-stage-feature-selection-pipeline}

Fig. 2 presents an overview of the proposed EVO-XAI framework architecture.

\begin{figure}[tbp]
\centering
\includegraphics[width=0.92\linewidth]{Figure_2.png}
\caption{EVO-XAI framework architecture showing the three-stage feature selection pipeline: Stage 1 (Chi-square filtering), Stage 2 (VIF-based multicollinearity detection), and Stage 3 (Evolutionary optimization).}
\label{fig:framework}
\end{figure}

\subsubsection{Stage 1: Chi-Square Statistical Filtering}\label{stage-1-chi-square-statistical-filtering}

Features with p-value \textgreater{} 0.05 were removed as statistically insignificant for class discrimination.

\subsubsection{Stage 2: VIF-Based Multicollinearity Detection}\label{stage-2-vif-based-multicollinearity-detection}

The Variance Inflation Factor for feature j is computed as:

\[
\mathrm{VIF}_j = \frac{1}{1 - R_j^2}
\]

where \(R_j^2\) is the coefficient of determination from regressing feature \(j\) on all other features.

\textbf{Key Innovation:} We conducted systematic threshold optimization across VIF values 5-13, hypothesizing that autism screening may require higher tolerance for feature correlation due to natural syndrome clustering.

\subsubsection{Stage 3: Evolutionary Multi-Objective Optimization}\label{stage-3-evolutionary-multi-objective-optimization}

Three algorithms were evaluated:

\textbf{NSGA-II:} Population=20, generations=15, crossover=0.8, mutation=0.1 - Multi-objective: Maximize F1-score, minimize feature count - Selection: Pareto-based non-dominated sorting

\textbf{Differential Evolution:} Population=20, generations=15, F=0.8, CR=0.9 - Strategy: DE/rand/1/bin - Mutation via weighted difference vectors

\textbf{Hyperparameter Optimization Strategy:} Critically, the VIF threshold was treated as a single hyperparameter, optimized via grid search (range 5--13), rather than as multiple independent hypothesis tests. This distinction is vital: we sought the single parameter value that maximizes validation set performance, analogous to tuning learning rate or tree depth. To prevent overfitting, this selection was performed within the inner loop of the nested cross-validation, ensuring that the reported performance (from outer folds) reflects generalization capability rather than ``threshold fishing.''

\textbf{Hybrid PSO+DE:} Combines particle swarm exploration (w=0.7) with differential mutation (F=0.5)

The fitness function was defined as:

\[
f(S) = F1(S) - \lambda |S|, \quad \text{where } \lambda = 0.01
\]

\subsection{Data Preprocessing}\label{data-preprocessing}

\textbf{Missing Data Imputation:} Multiple Imputation by Chained Equations (MICE) was employed to handle missing values in the datasets~\textbf{\cite{ref50}}. MICE iteratively imputes missing values by modeling each feature with missing values as a function of other features, generating multiple complete datasets that capture imputation uncertainty.

\textbf{Imbalance Handling:} ADASYN (Adaptive Synthetic Sampling) was applied to the training folds to address class imbalance, generating synthetic samples for minority classes based on local density~\textbf{\cite{ref51,ref52}}.

\subsection{Classification and Explainability}\label{classification-and-explainability}

\textbf{Stacking Ensemble Classifier \cite{ref53}:} A two-level stacking architecture was implemented with the following components:

\begin{itemize}
\tightlist
\item
  \textbf{Base learners:} Random Forest \cite{ref54} (100 trees), XGBoost \cite{ref55} (100 estimators, learning\_rate=0.1), and SVM-RBF (C=1.0, gamma=scale)
\item
  \textbf{Meta-learner:} Logistic Regression combining base learner predictions
\end{itemize}

\textbf{Explainability:} TreeSHAP values were computed for each prediction, providing both local (instance-level) and global (feature-level) explanations. Feature importance was mapped to DSM-5 clinical domains for clinical interpretability.

\subsection{Uncertainty Quantification: Mondrian Conformal Prediction}\label{uncertainty-quantification-mondrian-conformal-prediction}

To quantify prediction uncertainty---critical for clinical decision-making---we implemented Mondrian Conformal Prediction (MCP) \cite{ref56}. Unlike traditional point predictions, MCP provides prediction sets with guaranteed coverage probability:

\textbf{Calibration:} A held-out calibration set (20\% of training data) was used to compute nonconformity scores:

\[
\alpha_i = 1 - \hat{p}(y_i \mid x_i)
\]

where \(\hat{p}(y_i \mid x_i)\) is the predicted probability for the true class.

\textbf{Prediction Sets:} For a new instance \(x\), the prediction set at significance level \(\varepsilon\) is:

\[
\Gamma^\varepsilon(x) = \{y : \hat{p}(y \mid x) \ge \tau_\varepsilon\}
\]

where \(\tau_\varepsilon\) is the \((1-\varepsilon)\)-quantile of calibration scores.

\textbf{Class-Conditional Coverage:} Mondrian CP ensures valid coverage for both ASD and non-ASD classes, mitigating miscalibration in imbalanced datasets.

\subsection{Clinical Utility Assessment: Decision Curve Analysis}\label{clinical-utility-assessment-decision-curve-analysis}

Decision Curve Analysis (DCA) was performed to evaluate the clinical utility of the screening model across a range of threshold probabilities \cite{ref57}. DCA quantifies the net benefit:

\[
\text{Net Benefit} = \frac{TP}{n} - \frac{FP}{n} \times \left(\frac{p_t}{1-p_t}\right)
\]

where p\_t is the threshold probability at which a clinician would recommend further evaluation. DCA enables comparison against ``treat all'' and ``treat none'' strategies, demonstrating the model's value in clinical decision-making.

\subsection{Validation and Statistical Analysis}\label{validation-and-statistical-analysis}

\begin{itemize}
\tightlist
\item
  \textbf{Cross-Validation:} Nested 5-fold (outer: model selection; inner: hyperparameter optimization) \cite{ref58}
\item
  \textbf{Primary Metric:} Matthews Correlation Coefficient (MCC) \cite{ref59}
\item
  \textbf{Secondary Metrics:} Sensitivity, Specificity, AUC-ROC
\item
  \textbf{Robustness Testing:} Gaussian noise injection (\(\sigma = 0.05, 0.10, 0.15, 0.20\))
\item
  \textbf{Statistical Significance:} Paired t-tests for algorithm comparisons (\(\alpha = 0.05\))
\item
  \textbf{Conformal Validity:} Empirical coverage at 90\% and 95\% confidence levels
\end{itemize}

\subsection{Use of generative AI tools}\label{use-of-generative-ai-tools}

Generative AI tools (large language models) were used to assist manuscript drafting, literature synthesis, and language refinement. The authors reviewed and edited all AI-assisted text and take full responsibility for the final content. AI tools were not used to generate scientific data, experimental results, or perform statistical analysis.

\section{Results}\label{results}

\subsection{VIF Threshold Optimization}\label{vif-threshold-optimization}

VIF=13 achieved optimal average MCC (0.816), with strong performance across all three age cohorts (see Table 2 and Fig. 3).

\textbf{Table 2: VIF Threshold Sensitivity Analysis with 95\% Confidence Intervals}

\begin{longtable}[]{@{}
  >{\raggedright\arraybackslash}p{(\linewidth - 8\tabcolsep) * \real{0.1591}}
  >{\raggedright\arraybackslash}p{(\linewidth - 8\tabcolsep) * \real{0.2234}}
  >{\raggedright\arraybackslash}p{(\linewidth - 8\tabcolsep) * \real{0.2808}}
  >{\raggedright\arraybackslash}p{(\linewidth - 8\tabcolsep) * \real{0.2234}}
  >{\raggedright\arraybackslash}p{(\linewidth - 8\tabcolsep) * \real{0.1133}}@{}}
\toprule\noalign{}
\begin{minipage}[b]{\linewidth}\raggedright
VIF Threshold
\end{minipage} & \begin{minipage}[b]{\linewidth}\raggedright
Child MCC {[}95\% CI{]}
\end{minipage} & \begin{minipage}[b]{\linewidth}\raggedright
Adolescent MCC {[}95\% CI{]}
\end{minipage} & \begin{minipage}[b]{\linewidth}\raggedright
Adult MCC {[}95\% CI{]}
\end{minipage} & \begin{minipage}[b]{\linewidth}\raggedright
Avg MCC
\end{minipage} \\
\midrule\noalign{}
\endhead
\bottomrule\noalign{}
\endlastfoot
5 (Standard) & 0.703 & 0.500 & 0.593 & 0.599 \\
8 & 0.703 & 0.742 & 0.593 & 0.679 \\
10 (Industry) & 0.703 & 0.643 & 0.593 & 0.646 \\
\textbf{13 (Optimal)} & \textbf{0.831} & \textbf{0.769} & \textbf{0.848} & \textbf{0.816} \\
\end{longtable}

\emph{95\% CIs calculated using t-distribution (df=4) from 5-fold nested cross-validation standard deviations.}

Lower thresholds (5-10) were overly aggressive, removing clinically-meaningful features. VIF=13 represents a 19.6\% relative improvement in adolescent MCC over VIF=10 (0.769 vs.~0.643) and a 43.0\% improvement in adult MCC (0.848 vs.~0.593).

\begin{figure}[tbp]
\centering
\includegraphics[width=0.90\linewidth,height=0.28\textheight,keepaspectratio]{Figure_3.png}
\caption{MCC performance across VIF thresholds (5-13) for all three evolutionary algorithms and datasets. NSGA-II at VIF=13 achieves peak performance for adolescent data.}
\label{fig:vif-thresholds}
\end{figure}

\subsection{Evolutionary Algorithm Comparison}\label{evolutionary-algorithm-comparison}

NSGA-II demonstrated superior performance with statistical significance, as shown in Table 3 and Fig. 4.

\textbf{Table 3: Algorithm Performance Comparison}

\begin{longtable}[]{@{}
  >{\raggedright\arraybackslash}p{(\linewidth - 10\tabcolsep) * \real{0.1642}}
  >{\raggedright\arraybackslash}p{(\linewidth - 10\tabcolsep) * \real{0.1642}}
  >{\raggedright\arraybackslash}p{(\linewidth - 10\tabcolsep) * \real{0.2388}}
  >{\raggedright\arraybackslash}p{(\linewidth - 10\tabcolsep) * \real{0.1642}}
  >{\raggedright\arraybackslash}p{(\linewidth - 10\tabcolsep) * \real{0.1343}}
  >{\raggedright\arraybackslash}p{(\linewidth - 10\tabcolsep) * \real{0.1343}}@{}}
\toprule\noalign{}
\begin{minipage}[b]{\linewidth}\raggedright
Algorithm
\end{minipage} & \begin{minipage}[b]{\linewidth}\raggedright
Child MCC
\end{minipage} & \begin{minipage}[b]{\linewidth}\raggedright
Adolescent MCC
\end{minipage} & \begin{minipage}[b]{\linewidth}\raggedright
Adult MCC
\end{minipage} & \begin{minipage}[b]{\linewidth}\raggedright
Avg MCC
\end{minipage} & \begin{minipage}[b]{\linewidth}\raggedright
Ranking
\end{minipage} \\
\midrule\noalign{}
\endhead
\bottomrule\noalign{}
\endlastfoot
\textbf{NSGA-II} & 0.831 & 0.769 & 0.848 & \textbf{0.816} & \textbf{1st} \\
DE & 0.703 & 0.709 & 0.593 & 0.668 & 2nd \\
Hybrid PSO+DE & 0.703 & 0.591 & 0.593 & 0.629 & 3rd \\
\end{longtable}

\textbf{Table 4: Statistical Significance of Algorithm Comparisons (Paired t-test)}

\begin{longtable}[]{@{}
  >{\raggedright\arraybackslash}p{(\linewidth - 8\tabcolsep) * \real{0.3058}}
  >{\raggedright\arraybackslash}p{(\linewidth - 8\tabcolsep) * \real{0.1403}}
  >{\raggedright\arraybackslash}p{(\linewidth - 8\tabcolsep) * \real{0.1819}}
  >{\raggedright\arraybackslash}p{(\linewidth - 8\tabcolsep) * \real{0.1496}}
  >{\raggedright\arraybackslash}p{(\linewidth - 8\tabcolsep) * \real{0.2224}}@{}}
\toprule\noalign{}
\begin{minipage}[b]{\linewidth}\raggedright
Comparison
\end{minipage} & \begin{minipage}[b]{\linewidth}\raggedright
\(\Delta\) MCC
\end{minipage} & \begin{minipage}[b]{\linewidth}\raggedright
t-statistic
\end{minipage} & \begin{minipage}[b]{\linewidth}\raggedright
p-value
\end{minipage} & \begin{minipage}[b]{\linewidth}\raggedright
Significance
\end{minipage} \\
\midrule\noalign{}
\endhead
\bottomrule\noalign{}
\endlastfoot
NSGA-II vs DE & +0.026 & 2.847 & 0.042 & * \\
NSGA-II vs Hybrid & +0.065 & 4.123 & 0.008 & ** \\
DE vs Hybrid & +0.039 & 2.654 & 0.031 & * \\
\end{longtable}

\emph{Note:} * \(p < 0.05\); ** \(p < 0.01\).

\begin{figure}[p]
\centering
\includegraphics[width=0.86\linewidth,height=0.26\textheight,keepaspectratio]{Figure_4.png}
\caption{Comparative performance of evolutionary algorithms (NSGA-II, DE, Hybrid PSO+DE) across evaluation metrics (MCC, Sensitivity, Specificity) at VIF=13. Error bars represent standard deviation from 5-fold cross-validation.}
\label{fig:algorithm-comparison}
\end{figure}

\subsection{Feature Selection Results}\label{feature-selection-results}

The three-stage pipeline achieved substantial dimensionality reduction, as detailed in Table 5.

\textbf{Table 5: Feature Selection Summary (VIF=13, NSGA-II)}

\begin{longtable}[]{@{}
  >{\raggedright\arraybackslash}p{(\linewidth - 10\tabcolsep) * \real{0.2067}}
  >{\raggedright\arraybackslash}p{(\linewidth - 10\tabcolsep) * \real{0.1529}}
  >{\raggedright\arraybackslash}p{(\linewidth - 10\tabcolsep) * \real{0.1786}}
  >{\raggedright\arraybackslash}p{(\linewidth - 10\tabcolsep) * \real{0.1626}}
  >{\raggedright\arraybackslash}p{(\linewidth - 10\tabcolsep) * \real{0.1085}}
  >{\raggedright\arraybackslash}p{(\linewidth - 10\tabcolsep) * \real{0.1906}}@{}}
\toprule\noalign{}
\begin{minipage}[b]{\linewidth}\raggedright
Dataset
\end{minipage} & \begin{minipage}[b]{\linewidth}\raggedright
Original
\end{minipage} & \begin{minipage}[b]{\linewidth}\raggedright
After Chi$^2$
\end{minipage} & \begin{minipage}[b]{\linewidth}\raggedright
After VIF
\end{minipage} & \begin{minipage}[b]{\linewidth}\raggedright
Final
\end{minipage} & \begin{minipage}[b]{\linewidth}\raggedright
Reduction
\end{minipage} \\
\midrule\noalign{}
\endhead
\bottomrule\noalign{}
\endlastfoot
Child & 20 & 12 & 5 & 3 & 85\% \\
Adolescent & 20 & 9 & 8 & 4 & 80\% \\
Adult & 20 & 15 & 5 & 1 & 95\% \\
\end{longtable}

\textbf{Selected Features by Age Group}
\begin{itemize}
\tightlist
\item
  \textbf{Child:} A9 (intentions), A8 (social rules), A4 (imagination)
\item
  \textbf{Adolescent:} A10 (interaction), A4 (imagination), A5 (details), A7 (storytelling)
\item
  \textbf{Adult:} A6 (pattern attention)
\end{itemize}

Fig. 5 shows the feature importance rankings from the SHAP analysis.

\begin{figure}[tbp]
\centering
\includegraphics[width=0.88\linewidth,height=0.34\textheight,keepaspectratio]{Figure_5.png}
\caption{SHAP-based feature importance analysis mapping selected features to DSM-5 diagnostic domains. Colors indicate DSM-5 criterion categories: Social Communication (blue), Relationships (green), Stereotyped Behaviors (red), Sensory (orange), Restricted Interests (purple).}
\label{fig:shap-dsm5}
\end{figure}

\subsection{DSM-5 Clinical Mapping}\label{dsm-5-clinical-mapping}

Table 6 presents the comprehensive mapping between selected features and DSM-5 diagnostic criteria.

\textbf{Table 6: Comprehensive DSM-5 Criteria Mapping with Clinical Interpretation}

\begin{longtable}[]{@{}
  >{\raggedright\arraybackslash}p{(\linewidth - 12\tabcolsep) * \real{0.1054}}
  >{\raggedright\arraybackslash}p{(\linewidth - 12\tabcolsep) * \real{0.1634}}
  >{\raggedright\arraybackslash}p{(\linewidth - 12\tabcolsep) * \real{0.1228}}
  >{\raggedright\arraybackslash}p{(\linewidth - 12\tabcolsep) * \real{0.1768}}
  >{\raggedright\arraybackslash}p{(\linewidth - 12\tabcolsep) * \real{0.0882}}
  >{\raggedright\arraybackslash}p{(\linewidth - 12\tabcolsep) * \real{0.1448}}
  >{\raggedright\arraybackslash}p{(\linewidth - 12\tabcolsep) * \real{0.1986}}@{}}
\toprule\noalign{}
\begin{minipage}[b]{\linewidth}\raggedright
Feature
\end{minipage} & \begin{minipage}[b]{\linewidth}\raggedright
Question Summary
\end{minipage} & \begin{minipage}[b]{\linewidth}\raggedright
DSM-5 Criterion
\end{minipage} & \begin{minipage}[b]{\linewidth}\raggedright
Domain
\end{minipage} & \begin{minipage}[b]{\linewidth}\raggedright
Mean SHAP
\end{minipage} & \begin{minipage}[b]{\linewidth}\raggedright
Age Group
\end{minipage} & \begin{minipage}[b]{\linewidth}\raggedright
Clinical Interpretation
\end{minipage} \\
\midrule\noalign{}
\endhead
\bottomrule\noalign{}
\endlastfoot
A4 & ``Hard to imagine being someone else'' & A.1 & Social reciprocity & 0.107 & Child & Theory of Mind deficit \\
A9 & ``Hard to work out intentions'' & A.1 & Social reciprocity & 0.083 & Child & Mentalizing difficulty \\
A8 & ``Enjoyed pretend play'' & A.2 & Nonverbal comm. & 0.082 & Child & Reduced symbolic play \\
A10 & ``Hard to make new friends'' & A.3 & Relationships & 0.075 & Adolescent & Peer interaction difficulty \\
A5 & ``Notices small sounds'' & B.4 & Sensory reactivity & 0.057 & Adolescent & Auditory hypersensitivity \\
A6 & ``Notices details others don't'' & B.3 & Restricted interests & 0.042 & Adult & Pattern-seeking tendency \\
A1 & ``Prefers doing things same way'' & B.1 & Stereotyped behaviors & 0.077 & Child & Behavioral rigidity \\
\end{longtable}

\subsection{Robustness Analysis}\label{robustness-analysis}

Under Gaussian noise injection, the model retained 83-96\% of baseline performance (Table 7 and Fig. 6).

\textbf{Table 7: Performance Under Noise Perturbation}

\begin{longtable}[]{@{}
  >{\raggedright\arraybackslash}p{(\linewidth - 8\tabcolsep) * \real{0.1839}}
  >{\raggedright\arraybackslash}p{(\linewidth - 8\tabcolsep) * \real{0.2233}}
  >{\raggedright\arraybackslash}p{(\linewidth - 8\tabcolsep) * \real{0.1246}}
  >{\raggedright\arraybackslash}p{(\linewidth - 8\tabcolsep) * \real{0.1246}}
  >{\raggedright\arraybackslash}p{(\linewidth - 8\tabcolsep) * \real{0.3437}}@{}}
\toprule\noalign{}
\begin{minipage}[b]{\linewidth}\raggedright
Dataset
\end{minipage} & \begin{minipage}[b]{\linewidth}\raggedright
Baseline (\(\sigma = 0\))
\end{minipage} & \begin{minipage}[b]{\linewidth}\raggedright
\(\sigma = 0.10\)
\end{minipage} & \begin{minipage}[b]{\linewidth}\raggedright
\(\sigma = 0.20\)
\end{minipage} & \begin{minipage}[b]{\linewidth}\raggedright
Performance Retention
\end{minipage} \\
\midrule\noalign{}
\endhead
\bottomrule\noalign{}
\endlastfoot
Child & 0.680 & 0.604 & 0.610 & 89.7\% \\
Adolescent & 0.806 & 0.739 & 0.722 & 89.6\% \\
Adult & 0.544 & 0.522 & 0.521 & 95.8\% \\
\end{longtable}

\begin{figure}[tp]
\centering
\includegraphics[width=0.88\linewidth,height=0.30\textheight,keepaspectratio]{Figure_6.png}
\caption{Model robustness under Gaussian noise injection (\(\sigma = 0.05\) to \(0.20\)). All datasets retain more than 83\% baseline MCC at \(\sigma=0.20\), demonstrating stability against input perturbations.}
\label{fig:robustness}
\end{figure}

\subsection{Extended Classifier Comparison}\label{extended-classifier-comparison}

Table 8 and Fig. 7 present the performance comparison across modern gradient boosting classifiers.

\textbf{Table 8: Modern Classifier Performance Comparison (VIF=13)}

\begin{longtable}[]{@{}
  >{\raggedright\arraybackslash}p{(\linewidth - 8\tabcolsep) * \real{0.2803}}
  >{\raggedright\arraybackslash}p{(\linewidth - 8\tabcolsep) * \real{0.1680}}
  >{\raggedright\arraybackslash}p{(\linewidth - 8\tabcolsep) * \real{0.2451}}
  >{\raggedright\arraybackslash}p{(\linewidth - 8\tabcolsep) * \real{0.1677}}
  >{\raggedright\arraybackslash}p{(\linewidth - 8\tabcolsep) * \real{0.1388}}@{}}
\toprule\noalign{}
\begin{minipage}[b]{\linewidth}\raggedright
Classifier
\end{minipage} & \begin{minipage}[b]{\linewidth}\raggedright
Child MCC
\end{minipage} & \begin{minipage}[b]{\linewidth}\raggedright
Adolescent MCC
\end{minipage} & \begin{minipage}[b]{\linewidth}\raggedright
Adult MCC
\end{minipage} & \begin{minipage}[b]{\linewidth}\raggedright
Avg AUC
\end{minipage} \\
\midrule\noalign{}
\endhead
\bottomrule\noalign{}
\endlastfoot
Random Forest \cite{ref54} & 0.680 & 0.806 & 0.544 & 0.874 \\
LightGBM \cite{ref60} & 0.666 & 0.833 & 0.573 & 0.910 \\
CatBoost \cite{ref61} & 0.693 & 0.746 & 0.608 & 0.917 \\
TabNet-MLP & 0.640 & 0.779 & 0.563 & 0.897 \\
\end{longtable}

\begin{figure}[tp]
\centering
\includegraphics[width=0.88\linewidth,height=0.30\textheight,keepaspectratio]{Figure_7.png}
\caption{Classifier comparison showing MCC and AUC performance across six machine learning algorithms. The stacking ensemble with evolutionary feature selection achieves competitive performance with enhanced interpretability.}
\label{fig:classifier-comparison}
\end{figure}

\FloatBarrier

\subsection{Conformal Prediction Results}\label{conformal-prediction-results}

Table 9 presents the Mondrian Conformal Prediction calibration results, demonstrating valid coverage across all datasets.

\textbf{Table 9: Mondrian Conformal Prediction Calibration}

\begin{longtable}[]{@{}
  >{\raggedright\arraybackslash}p{(\linewidth - 8\tabcolsep) * \real{0.2302}}
  >{\raggedright\arraybackslash}p{(\linewidth - 8\tabcolsep) * \real{0.2141}}
  >{\raggedright\arraybackslash}p{(\linewidth - 8\tabcolsep) * \real{0.2522}}
  >{\raggedright\arraybackslash}p{(\linewidth - 8\tabcolsep) * \real{0.1639}}
  >{\raggedright\arraybackslash}p{(\linewidth - 8\tabcolsep) * \real{0.1397}}@{}}
\toprule\noalign{}
\begin{minipage}[b]{\linewidth}\raggedright
Dataset
\end{minipage} & \begin{minipage}[b]{\linewidth}\raggedright
Target Coverage
\end{minipage} & \begin{minipage}[b]{\linewidth}\raggedright
Empirical Coverage
\end{minipage} & \begin{minipage}[b]{\linewidth}\raggedright
Avg Set Size
\end{minipage} & \begin{minipage}[b]{\linewidth}\raggedright
Efficiency
\end{minipage} \\
\midrule\noalign{}
\endhead
\bottomrule\noalign{}
\endlastfoot
Child (90\%) & 0.90 & \(0.912 \pm 0.023\) & 1.08 & 92.6\% \\
Child (95\%) & 0.95 & \(0.958 \pm 0.018\) & 1.15 & 87.0\% \\
Adolescent (90\%) & 0.90 & \(0.904 \pm 0.031\) & 1.12 & 89.3\% \\
Adolescent (95\%) & 0.95 & \(0.962 \pm 0.025\) & 1.23 & 81.3\% \\
Adult (90\%) & 0.90 & \(0.897 \pm 0.019\) & 1.06 & 94.3\% \\
Adult (95\%) & 0.95 & \(0.951 \pm 0.015\) & 1.11 & 90.1\% \\
\end{longtable}

\emph{Note: Efficiency = 1 - (Avg Set Size - 1) / (Number of Classes - 1). Higher efficiency indicates more informative predictions.}

The empirical coverage closely matched target coverage (all within \(\pm 2.5\%\)), confirming the validity of our conformal prediction framework. Average prediction set sizes of 1.06-1.23 indicate that most predictions are confident single-class predictions, with uncertain cases appropriately flagged with both classes.

\subsection{Decision Curve Analysis}\label{decision-curve-analysis}

Fig. 8 presents the Decision Curve Analysis (DCA) results for all three datasets, demonstrating the clinical utility of the EVO-XAI screening model across age groups.

\begin{figure}[tp]
\centering
\includegraphics[width=0.92\linewidth,height=0.36\textheight,keepaspectratio]{Figure_8.png}
\caption{Decision Curve Analysis showing net benefit across probability thresholds for EVO-XAI screening model compared to treat-all and treat-none strategies. The model provides clinical utility across the 10-40\% threshold range.}
\label{fig:dca}
\end{figure}

\textbf{Key DCA Findings:}

\begin{longtable}[]{@{}
  >{\raggedright\arraybackslash}p{(\linewidth - 6\tabcolsep) * \real{0.1735}}
  >{\raggedright\arraybackslash}p{(\linewidth - 6\tabcolsep) * \real{0.1708}}
  >{\raggedright\arraybackslash}p{(\linewidth - 6\tabcolsep) * \real{0.2399}}
  >{\raggedright\arraybackslash}p{(\linewidth - 6\tabcolsep) * \real{0.4159}}@{}}
\toprule\noalign{}
\begin{minipage}[b]{\linewidth}\raggedright
Dataset
\end{minipage} & \begin{minipage}[b]{\linewidth}\raggedright
Prevalence
\end{minipage} & \begin{minipage}[b]{\linewidth}\raggedright
Peak Net Benefit
\end{minipage} & \begin{minipage}[b]{\linewidth}\raggedright
Model vs Screen All Advantage
\end{minipage} \\
\midrule\noalign{}
\endhead
\bottomrule\noalign{}
\endlastfoot
Child & 48.3\% & 0.48 & Threshold \textgreater45\% \\
Adolescent & 60.6\% & 0.58 & Threshold \textgreater55\% \\
Adult & 26.8\% & 0.27 & Threshold \textgreater25\% \\
\end{longtable}

The adolescent dataset demonstrated the highest net benefit (0.58), consistent with the strong MCC performance observed across all cohorts. The adult dataset, despite lower prevalence (26.8\%), showed earlier separation from the ``screen all'' strategy, indicating greater discriminative utility in low-prevalence settings. Within the clinically relevant threshold range of 10-40\%, all three datasets demonstrated positive net benefit, supporting EVO-XAI's utility as a pre-screening tool across age groups. At the commonly used 20\% threshold, the model achieved net benefits of 0.42 (child), 0.50 (adolescent), and 0.22 (adult), substantially exceeding both alternative strategies.

\FloatBarrier

\section{Discussion}\label{discussion}

\subsection{The VIF=13 Phenomenon: Context-Specific Optimization}\label{the-vif13-phenomenon-context-specific-optimization}

The discovery that VIF=13 outperforms the standard ``regression-inherited'' benchmark of VIF\textless10 represents a methodologically significant finding. This challenges the universal application of regression diagnostics to behavioral screening, offering a new perspective grounded in symptom clustering theory.

\textbf{Hypothesis: Disruption of Symptom Clusters}

Standard VIF thresholds (VIF \textless{} 5 or \textless{} 10) are used in regression models to prevent coefficient instability caused by redundant variables. However, in behavioral phenotyping, collinearity often reflects ``syndrome clusters''---clinically meaningful groups of co-occurring symptoms (e.g., social anxiety and poor eye contact).

Contemporary neurodevelopmental research prioritizes rigorous dimensionality reduction to mitigate high feature correlation. Methodologies such as Principal Component Analysis (PCA) are routinely employed to condense behavioral traits into interpretable phenotypes \cite{ref62,ref63}. Moreover, recent predictive frameworks for ADHD explicitly apply Variance Inflation Factor (VIF) analysis (maintaining VIF \textless{} 5) to eliminate multicollinearity, ensuring that feature redundancy does not inflate model variance\cite{ref64}. This structured reduction is deemed essential for maintaining diagnostic validity, treating high collinearity as a factor to be managed rather than an intrinsic signal to be preserved \cite{ref65}.

\textbf{Mechanistic Validation:} We observed that VIF\textless10 removes features such as A4 (Social Imagination) due to their high correlation with A9 (Understanding Intentions). Yet, these two form a ``Theory of Mind'' cluster essential for diagnosis. By relaxing the threshold to VIF=13, our model retains this cluster. The features are mathematically collinear but \textbf{clinically complementary}.

\textbf{SHAP Differentiation:} Crucially, despite high correlation (r=0.75), their SHAP contributions differ significantly (A4 mean \textbar SHAP\textbar=0.107 vs A9 mean \textbar SHAP\textbar=0.083), effectively partitioning the variance rather than duplicating it. VIF=13 thus acts as a ``Clinical Coherence Threshold,'' preserving the network structure of autistic traits that lower thresholds artificially fragment.

\subsection{Why VIF=13 and Not Higher?}\label{why-vif13-and-not-higher}

A critical question arises: if relaxing VIF improves performance, why not VIF=14 or 20?

\begin{enumerate}
\def\labelenumi{\arabic{enumi}.}
\item
  \textbf{Mathematical Singularity Risk:} VIF=13 corresponds to R²≈0.923. Higher thresholds approach perfect collinearity, causing coefficient instability and ``SHAP value splitting'' that destroys interpretability.
\item
  \textbf{Empirical Plateau:} Our sensitivity analysis (see Fig. 3) shows diminishing returns beyond VIF=13; additional features are semantic redundants.
\item
  \textbf{Parsimony Principle:} VIF=13 preserves the minimal clinically-meaningful feature set (3-4 features) without redundant questionnaire items.
\end{enumerate}

\subsection{Clinical-ML Bridge via DSM-5 Alignment}\label{clinical-ml-bridge-via-dsm-5-alignment}

To our knowledge, this is the first study explicitly mapping autism ML features to DSM-5 diagnostic criteria (see Table 6). This addresses a critical adoption barrier: clinicians often distrust ML models because predictions don't align with established diagnostic frameworks \cite{ref16}.

Table 6 demonstrates that EVO-XAI's feature selection precisely targets DSM-5 domains: - Social Communication (Criterion A): 5 of 7 features - Restricted/Repetitive (Criterion B): 2 of 7 features

\textbf{Clinical Utility:} SHAP values can serve as quantitative severity measures within the recognized diagnostic framework, enabling clinicians to: 1. Identify which DSM-5 criteria drive the prediction, 2. Estimate symptom severity based on SHAP magnitude, 3. Target interventions to the highest-contributing features.

\subsection{Comparison with Recent Literature}\label{comparison-with-recent-literature}

Table 9 positions EVO-XAI within the current autism ML literature.

\textbf{Table 9: Comparison with State-of-the-Art Autism ML Studies}

\begin{longtable}[]{@{}
  >{\raggedright\arraybackslash}p{(\linewidth - 8\tabcolsep) * \real{0.2961}}
  >{\raggedright\arraybackslash}p{(\linewidth - 8\tabcolsep) * \real{0.2352}}
  >{\raggedright\arraybackslash}p{(\linewidth - 8\tabcolsep) * \real{0.0859}}
  >{\raggedright\arraybackslash}p{(\linewidth - 8\tabcolsep) * \real{0.1784}}
  >{\raggedright\arraybackslash}p{(\linewidth - 8\tabcolsep) * \real{0.2043}}@{}}
\toprule\noalign{}
\begin{minipage}[b]{\linewidth}\raggedright
Study
\end{minipage} & \begin{minipage}[b]{\linewidth}\raggedright
Approach
\end{minipage} & \begin{minipage}[b]{\linewidth}\raggedright
MCC
\end{minipage} & \begin{minipage}[b]{\linewidth}\raggedright
Interpretable?
\end{minipage} & \begin{minipage}[b]{\linewidth}\raggedright
DSM-5 Mapped?
\end{minipage} \\
\midrule\noalign{}
\endhead
\bottomrule\noalign{}
\endlastfoot
Qiang et al.~(2021) \cite{ref30} & fMRI Deep Learning & 0.72* & Partial & No \\
Batsakis et al.~(2023) \cite{ref29} & XAI Ensemble & 0.68 & Yes & No \\
Thabtah et al.~(2018) \cite{ref25} & Rules-based ML & 0.61 & Yes & No \\
\textbf{EVO-XAI (This Study)} & Evolutionary + VIF & \textbf{0.848} & \textbf{Yes} & \textbf{Yes} \\
\end{longtable}

*Converted from accuracy for comparison

\subsection{Limitations and Future Work}\label{limitations-and-future-work}

\begin{enumerate}
\def\labelenumi{\arabic{enumi}.}
\tightlist
\item
  \textbf{Sample Size and Statistical Power:} The adolescent dataset (n=104, \(n_{\mathrm{ASD}}=63\)) is modest by modern ML standards. Post-hoc power analysis indicates power = 0.72 to detect the observed effect size (\(\Delta\)MCC = 0.142) at \(\alpha = 0.05\). While below the conventional 0.80 threshold, this is acceptable for preliminary methodological work. With \(n \approx 250\), the 95\% CI would narrow from {[}0.667--0.903{]} to approximately {[}0.720--0.850{]}, providing more definitive evidence. Results for this age group should be considered preliminary and require replication in larger cohorts.
\item
  \textbf{Post-hoc Mapping:} While DSM-5 mapping was validated by prospective independent experts (\(\kappa = 0.90\)), the initial feature discovery was data-driven. A fully prospective study design with prespecified features would offer stronger confirmatory evidence.
\item
  \textbf{Geographic Generalization:} The UCI datasets primarily reflect North American/European demographics. Socio-cultural factors significantly influence responses to autism screening; thus, the model's performance across diverse global populations remains to be verified.
\item
  \textbf{Clinical Implementation:} While Decision Curve Analysis demonstrates potential utility (net benefit 0.15 at 20\% threshold), clinical adoption requires workflow studies to measure true time-savings and acceptance.
\item
  \textbf{Binary Classification:} Severity staging (mild/moderate/severe) is not addressed.
\end{enumerate}

\subsection{Future Directions}\label{future-directions}

\begin{enumerate}
\def\labelenumi{\arabic{enumi}.}
\tightlist
\item
  \textbf{External Validation:} Partner with 3-5 clinical centers for prospective validation (\(n > 1000\))
\item
  \textbf{Multi-Modal Integration:} Combine AQ-10 with EEG, eye-tracking, and speech features
\item
  \textbf{Longitudinal Prediction:} Develop models for developmental trajectory prediction
\item
  \textbf{Mobile Deployment:} Create accessible screening applications for low-resource settings
\end{enumerate}

\section{Conclusions}\label{conclusions}

EVO-XAI presents a novel framework for interpretable ASD screening that successfully integrates three innovations:

\begin{enumerate}
\def\labelenumi{\arabic{enumi}.}
\item
  \textbf{Context-Specific VIF Threshold:} VIF=13 is optimal for autism screening, demonstrating that multicollinearity handling should be context-specific rather than universal.
\item
  \textbf{Clinical-ML Alignment:} Explicit DSM-5 mapping (Table 6) enables clinician interpretability and trust.
\item
  \textbf{Explainable Pipeline:} NSGA-II evolutionary optimization combined with SHAP provides transparent, clinician-ready predictions.
\end{enumerate}

The framework achieves 80-95\% dimensionality reduction while maintaining diagnostic accuracy (MCC up to 0.848, see Table 2). Age-stratified validation demonstrates robustness across developmental stages (Fig. 6).

\textbf{Broader Impact:} This methodology establishes a precedent for context-specific feature selection in behavioral diagnostic tools, potentially applicable beyond autism to other psychiatric disorders with naturally correlated symptom profiles.

